%! Unidad 3 — Estructura atómica y tabla periódica: ejercicios resueltos y ejercicios finales
%! Fuente: Fuentes/Quimica/Unidad 3 («Ejercicios resueltos», «Practica y calcula», «Piensa y resuelve»).
%! Enunciados adaptados con datos cambiados; soluciones recalculadas. Dificultad: \(\star\) baja, \(\star\star\) media, \(\star\star\star\) alta.
%! Convención para el electrón diferenciador: los orbitales se llenan de m = -l a m = +l, primero con s = +1/2.

\section{Ejercicios resueltos}\label{sec:ud3-resueltos}

Salvo que se indique otra cosa, usa \(h = \qty{6.626e-34}{\joule\s}\); \(c = \qty{3.00e8}{\m\per\s}\);
\(N_A = \qty{6.022e23}{\per\mol}\); \(R = \qty{1.097e7}{\per\m}\); \(E_n = -\qty{2.18e-18}{\joule}/n^2\) para el hidrógeno y
\(\qty{1}{\electronvolt} = \qty{1.602e-19}{\joule}\).

\begin{example}
    \textbf{¿Qué luz tiene más energía?} En el espectro visible, la luz amarilla tiene una longitud de onda de unos
    \qty{589}{\nm} y la violeta, de unos \qty{400}{\nm}. ¿Cuál de estas afirmaciones es correcta? a) La luz amarilla tiene más
    energía que la violeta. b) Las dos tienen la misma energía. c) La violeta tiene más energía que la amarilla. d) Las dos son más
    energéticas que los rayos X.

    \textbf{Solución.} \textbf{1. Datos}: \(\lambda_{\text{amarilla}} = \qty{589}{\nm}\); \(\lambda_{\text{violeta}} = \qty{400}{\nm}\).
    Se pide comparar la energía de sus fotones.

    \textbf{2. Cálculos.} Con la ecuación de Planck:
    \begin{gather*}
        E_{\text{amarilla}} = \frac{h c}{\lambda} = \frac{\num{6.626e-34} \cdot \num{3.00e8}}{\num{589e-9}} = \qty{3.37e-19}{\joule} \\
        E_{\text{violeta}} = \frac{\num{6.626e-34} \cdot \num{3.00e8}}{\num{400e-9}} = \qty{4.97e-19}{\joule}
    \end{gather*}
    La energía es inversamente proporcional a la longitud de onda: la luz violeta es más energética que la amarilla. Los rayos X,
    con \(\lambda \approx \qty{1}{\nm}\), tienen fotones de unos \qty{2e-16}{\joule}, cientos de veces más energéticos que los
    visibles. La respuesta correcta es la \textbf{c}.
\end{example}

\begin{example}
    \textbf{Números cuánticos permitidos.} Dados los conjuntos de números cuánticos \((2, 2, 0, +\frac{1}{2})\),
    \((4, 1, 0, +\frac{1}{2})\), \((3, 1, 0, 2)\), \((3, 2, -2, -\frac{1}{2})\) y \((3, 1, 2, +\frac{1}{2})\): a) ¿cuáles están
    permitidos y cuáles no? b) ¿En qué orbital está el electrón en los permitidos?

    \textbf{Solución.} Recordamos las reglas: \(l = 0, 1\dots, n - 1\); \(m\) va de \(-l\) a \(+l\); \(m_s = \pm\frac{1}{2}\).

    a) \((2, 2, 0, +\frac{1}{2})\): \textbf{no}, \(l\) no puede ser igual a \(n\). \((4, 1, 0, +\frac{1}{2})\): \textbf{sí}.
    \((3, 1, 0, 2)\): \textbf{no}, \(m_s\) solo puede valer \(\pm\frac{1}{2}\). \((3, 2, -2, -\frac{1}{2})\): \textbf{sí}.
    \((3, 1, 2, +\frac{1}{2})\): \textbf{no}, con \(l = 1\), \(m\) solo puede ser \(-1\), 0 o \(+1\).

    b) \(n\) indica el nivel y \(l\), el tipo de orbital: \((4, 1, 0, +\frac{1}{2})\) es un electrón \textbf{4p}, y
    \((3, 2, -2, -\frac{1}{2})\), un electrón \textbf{3d}.
\end{example}

\begin{example}
    \textbf{Configuración y propiedades periódicas.} Para los elementos A (\(Z = 17\)), B (\(Z = 35\)), C (\(Z = 19\)) y
    D (\(Z = 20\)): a) escribe su configuración electrónica; b) indica su periodo y su grupo; c) ¿cuál tiene mayor radio atómico?;
    d) ¿cuál tiene mayor energía de ionización?; e) ¿cuál tiene mayor afinidad electrónica? Razona las respuestas.

    \textbf{Solución.} a) Como son átomos neutros, tienen \(Z\) electrones. Con el diagrama de Möller:
    \begin{center}
        A: \conf{1s^2\,2s^2\,2p^6\,3s^2\,3p^5} \qquad B: \conf{1s^2\,2s^2\,2p^6\,3s^2\,3p^6\,4s^2\,3d^{10}\,4p^5} \\[2pt]
        C: \conf{1s^2\,2s^2\,2p^6\,3s^2\,3p^6\,4s^1} \qquad D: \conf{1s^2\,2s^2\,2p^6\,3s^2\,3p^6\,4s^2}
    \end{center}

    b) A: capa de valencia \conf{3s^2\,3p^5}, periodo 3, grupo 17 (es el cloro). B: \conf{4s^2\,4p^5}, periodo 4, grupo 17
    (bromo). C: \conf{4s^1}, periodo 4, grupo 1 (potasio). D: \conf{4s^2}, periodo 4, grupo 2 (calcio).

    c) El radio aumenta al bajar en un grupo (hay más capas) y disminuye hacia la derecha en un periodo (aumenta la carga nuclear
    efectiva). B, C y D están en el periodo 4, y C es el que está más a la izquierda: \textbf{C} tiene el mayor radio.

    d) La energía de ionización aumenta hacia la derecha y hacia arriba, al contrario que el radio: los electrones están más
    cerca del núcleo y más atraídos. El mayor valor corresponde a \textbf{A}, el que está más arriba y más a la derecha.

    e) La afinidad electrónica (la energía que se desprende al captar un electrón) también aumenta hacia arriba y hacia la
    derecha: \textbf{A} es el que tiene mayor tendencia a captar un electrón y completar su capa de valencia.
\end{example}

\begin{example}
    \textbf{Una raya del espectro del hidrógeno.} El electrón de un átomo de hidrógeno excitado pasa del nivel \(n = 4\) al
    \(n = 2\). Calcula la energía, la frecuencia y la longitud de onda del fotón emitido. ¿De qué color es esa raya? ¿A qué serie
    pertenece?

    \textbf{Solución.} \textbf{1. Energía}: es la diferencia entre los dos niveles:
    \[ \Delta E = E_4 - E_2 = -\frac{\num{2.18e-18}}{16} - \left( -\frac{\num{2.18e-18}}{4} \right) = \num{2.18e-18} \left( \frac{1}{4} - \frac{1}{16} \right) = \qty{4.09e-19}{\joule} \]
    \textbf{2. Frecuencia y longitud de onda}:
    \[ \nu = \frac{\Delta E}{h} = \frac{\num{4.09e-19}}{\num{6.626e-34}} = \qty{6.17e14}{\hertz} \qquad
    \lambda = \frac{c}{\nu} = \frac{\num{3.00e8}}{\num{6.17e14}} = \qty{4.86e-7}{\m} = \qty{486}{\nm} \]
    Es una radiación visible, de color \textbf{azul verdoso}. Como el salto termina en \(n = 2\), pertenece a la serie de
    \textbf{Balmer}. Se obtiene lo mismo con la fórmula de Rydberg: \(1/\lambda = R(1/4 - 1/16)\).
\end{example}

\begin{example}
    \textbf{Abundancia de los isótopos del cobre.} El cobre natural tiene dos isótopos, \ch{^{63}Cu} (\qty{62.930}{\uma}) y
    \ch{^{65}Cu} (\qty{64.928}{\uma}), y su masa atómica es \qty{63.546}{\uma}. Calcula la abundancia de cada isótopo y dibuja el
    espectro de masas.

    \textbf{Solución.} Llamamos \(x\) al porcentaje de \ch{^{63}Cu}; el de \ch{^{65}Cu} será \(100 - x\):
    \[ \num{63.546} = \frac{\num{62.930}\,x + \num{64.928}\,(100 - x)}{100} \;\Rightarrow\; \num{6354.6} = \num{6492.8} - \num{1.998}\,x
    \;\Rightarrow\; x = \num{69.2} \]
    Hay un \textbf{69,2\,\%} de \ch{^{63}Cu} y un \textbf{30,8\,\%} de \ch{^{65}Cu}. El espectro de masas tiene dos picos, en
    \(\approx 63\) y \(\approx 65\), y el primero es algo más del doble de alto que el segundo.
\end{example}

%========================================================================================
\section{Ejercicios de aplicación}\label{sec:ud3-aplicacion}
%========================================================================================

\subsection{Modelos atómicos}

\begin{exercise}
    \(\star\) ¿Qué masa y qué carga tienen un protón, un electrón y un neutrón? ¿Qué tienen en común estas tres partículas? ¿En
    qué se diferencian?
    \begin{solu}
        \textcolor{red}{Protón: \qty{1.673e-27}{\kilo\gram}, \(+\qty{1.602e-19}{\coulomb}\). Electrón: \qty{9.109e-31}{\kilo\gram},
        \(-\qty{1.602e-19}{\coulomb}\). Neutrón: \qty{1.675e-27}{\kilo\gram}, sin carga.} Las tres forman parte de todos los
        átomos (salvo el \ch{^{1}H}, sin neutrones). Se diferencian en la carga, en la masa (protón y neutrón tienen masas casi
        iguales, unas 1840 veces la del electrón) y en su localización: protones y neutrones en el núcleo, electrones en la corteza.
    \end{solu}
\end{exercise}

\begin{exercise}
    \(\star\) Dibuja en tu cuaderno el modelo atómico de Thomson y el de Rutherford e indica en cada dibujo dónde están la carga
    positiva, la carga negativa y la masa.
    \begin{solu}
        \textcolor{red}{Thomson}: una esfera con la carga positiva y la masa repartidas uniformemente y los electrones incrustados.
        \textcolor{red}{Rutherford}: un núcleo diminuto, con toda la carga positiva y casi toda la masa, y los electrones girando a
        su alrededor en un espacio casi vacío (figura~\ref{fig:ud3-modelos}).
    \end{solu}
\end{exercise}

\begin{exercise}
    \(\star\star\) ¿Cuáles son las diferencias principales entre el modelo atómico de Thomson y el de Rutherford? ¿Qué experimento
    obligó a cambiar de uno a otro?
    \begin{solu}
        \textcolor{red}{En el de Thomson el átomo es una esfera maciza con la carga positiva repartida; en el de Rutherford está
        casi vacío, la carga positiva y la masa se concentran en un núcleo pequeñísimo y los electrones giran en la corteza.} El
        cambio lo provocó \textcolor{red}{el experimento de la lámina de oro} (1910): algunas partículas \(\alpha\) se desviaban
        mucho o rebotaban, cosa imposible con el modelo de Thomson.
    \end{solu}
\end{exercise}

\begin{exercise}
    \(\star\star\) ¿Cuáles son las limitaciones del modelo atómico de Rutherford?
    \begin{solu}
        \textcolor{red}{Según la física clásica, el electrón, al girar, emitiría radiación, perdería energía y caería sobre el
        núcleo: el átomo no sería estable.} Además, \textcolor{red}{no explicaba los espectros discontinuos} de los átomos.
    \end{solu}
\end{exercise}

\begin{exercise}
    \(\star\) ¿Cuál es la diferencia entre una órbita y un orbital?
    \begin{solu}
        \textcolor{red}{Una órbita (modelo de Bohr) es una trayectoria circular definida por la que se mueve el electrón; un
        orbital (modelo mecanocuántico) es una región del espacio donde la probabilidad de encontrar al electrón es máxima}, sin
        trayectoria definida, porque no se pueden conocer a la vez su posición y su velocidad.
    \end{solu}
\end{exercise}

\begin{exercise}
    \(\star\) ¿A qué modelos atómicos se les llama «pudin de pasas» y «planetario»? ¿Qué aportó cada uno al conocimiento del átomo?
    \begin{solu}
        \textcolor{red}{«Pudin de pasas»: el de Thomson}; aportó que el átomo es divisible y contiene electrones.
        \textcolor{red}{«Planetario»: el de Rutherford}; aportó la existencia del núcleo y que el átomo está casi vacío.
    \end{solu}
\end{exercise}

\begin{exercise}
    \(\star\star\) Los modelos sobre la estructura del átomo empezaron a desarrollarse a comienzos del siglo XX. a) ¿A qué crees
    que se debe? b) Desde entonces no han dejado de aparecer nuevos modelos y aportaciones. ¿Por qué hubo tantos avances en tan
    poco tiempo?
    \begin{solu}
        a) \textcolor{red}{A los descubrimientos de finales del siglo XIX}: el electrón, los rayos X, la radiactividad y la
        espectroscopía mostraron que el átomo tenía estructura interna. b) \textcolor{red}{Cada modelo se ponía a prueba con nuevos
        experimentos, cada vez más precisos, y la física cuántica dio herramientas nuevas}; además, los científicos trabajaban en
        contacto y aprovechaban los resultados de los demás.
    \end{solu}
\end{exercise}

\begin{exercise}
    \(\star\star\) Haced en grupo un eje cronológico de la evolución de los modelos atómicos, con los científicos que intervinieron
    y su aportación, y exponedlo en clase.
    \begin{solu}
        \textcolor{red}{Respuesta abierta.} Por ejemplo: Demócrito (s. V a.\,C., idea de átomo); Dalton (1808, esferas
        indivisibles); Thomson (1897, electrón; 1904, modelo); Millikan (1909, carga del electrón); Rutherford (1911, núcleo;
        1919, protón); Bohr (1913, niveles de energía); Sommerfeld (1915, subniveles); De Broglie (1924, dualidad onda-partícula);
        Schrödinger y Heisenberg (1926-1927, modelo mecanocuántico); Chadwick (1932, neutrón).
    \end{solu}
\end{exercise}

\begin{exercise}
    \(\star\star\) El radio de un átomo de oro es de unos \qty{1.44e-10}{\m} y el de su núcleo, de unos \qty{7.3e-15}{\m}.
    Considerando ambos esféricos, ¿qué fracción del volumen del átomo ocupa el núcleo?
    \begin{solu}
        \(\dfrac{V_{\text{núcleo}}}{V_{\text{átomo}}} = \left( \dfrac{\num{7.3e-15}}{\num{1.44e-10}} \right)^3 =
        \textcolor{red}{\num{1.3e-13}}\): el núcleo ocupa una parte en casi diez billones del volumen del átomo.
    \end{solu}
\end{exercise}

\subsection{Masa atómica}

\begin{exercise}
    \(\star\) Indica el número másico, el número atómico y el número de protones, electrones y neutrones de estos átomos neutros:
    \ch{^{58}_{28}Ni}, \ch{^{87}_{38}Sr}, \ch{^{40}_{20}Ca}, \ch{^{17}_{8}O} y \ch{^{127}_{53}I}.
    \begin{solu}
        \begin{tabular}{@{}lccccc@{}}
            \toprule
            & \ch{^{58}_{28}Ni} & \ch{^{87}_{38}Sr} & \ch{^{40}_{20}Ca} & \ch{^{17}_{8}O} & \ch{^{127}_{53}I} \\
            \midrule
            \(A\) / \(Z\) & 58 / 28 & 87 / 38 & 40 / 20 & 17 / 8 & 127 / 53 \\[3pt]
            Protones = electrones & 28 & 38 & 20 & 8 & 53 \\[3pt]
            Neutrones & \textcolor{red}{30} & \textcolor{red}{49} & \textcolor{red}{20} & \textcolor{red}{9} & \textcolor{red}{74} \\
            \bottomrule
        \end{tabular}
    \end{solu}
\end{exercise}

\begin{exercise}
    \(\star\star\) Indica el número de protones, electrones y neutrones de estos iones: \ch{^{27}_{13}Al^3+}, \ch{^{133}_{55}Cs^+},
    \ch{^{24}_{12}Mg^2+}, \ch{^{34}_{16}S^2-} y \ch{^{81}_{35}Br^-}.
    \begin{solu}
        \ch{Al^3+}: \textcolor{red}{13 p, 10 e, 14 n}. \ch{Cs^+}: \textcolor{red}{55 p, 54 e, 78 n}. \ch{Mg^2+}:
        \textcolor{red}{12 p, 10 e, 12 n}. \ch{S^2-}: \textcolor{red}{16 p, 18 e, 18 n}. \ch{Br^-}: \textcolor{red}{35 p, 36 e, 46 n}.
    \end{solu}
\end{exercise}

\begin{exercise}
    \(\star\star\) Determina el número de protones, electrones y neutrones de estas parejas: a) \ch{^{31}_{15}P^3-} y
    \ch{^{31}_{15}P}; b) \ch{^{138}_{56}Ba^2+} y \ch{^{138}_{56}Ba}. ¿Qué cambia y qué no al formarse el ion?
    \begin{solu}
        a) \ch{P^3-}: \textcolor{red}{15 p, 18 e, 16 n}; P: \textcolor{red}{15 p, 15 e, 16 n}. b) \ch{Ba^2+}: \textcolor{red}{56 p,
        54 e, 82 n}; Ba: \textcolor{red}{56 p, 56 e, 82 n}. \textcolor{red}{Al formarse un ion solo cambia el número de electrones}:
        el núcleo no varía.
    \end{solu}
\end{exercise}

\begin{exercise}
    \(\star\star\) El cobre es un oligoelemento esencial: ayuda a mantener los vasos sanguíneos, los nervios y los huesos, y a
    formar los glóbulos rojos. Puede formar iones \ch{Cu^+} y \ch{Cu^2+}. Indica el número de protones, neutrones y electrones de
    \ch{^{63}_{29}Cu^+} y de \ch{^{63}_{29}Cu^2+}.
    \begin{solu}
        Los dos tienen \textcolor{red}{29 protones y 34 neutrones}; \ch{Cu^+} tiene \textcolor{red}{28 electrones} y \ch{Cu^2+},
        \textcolor{red}{27}.
    \end{solu}
\end{exercise}

\begin{exercise}
    \(\star\) El cobalto tiene 22 isótopos conocidos. El cobalto-60 se usa en radioterapia contra el cáncer. Sabiendo que para el
    cobalto \(Z = 27\), indica cuántos protones, neutrones y electrones tiene un átomo de \ch{^{60}Co}.
    \begin{solu}
        \textcolor{red}{27 protones, 33 neutrones y 27 electrones.}
    \end{solu}
\end{exercise}

\begin{exercise}
    \(\star\star\) El espectro de masas del magnesio presenta tres picos: \ch{^{24}Mg} (\qty{23.985}{\uma}; 78,99\,\%),
    \ch{^{25}Mg} (\qty{24.986}{\uma}; 10,00\,\%) y \ch{^{26}Mg} (\qty{25.983}{\uma}; 11,01\,\%). Represéntalo y calcula la masa
    atómica media del magnesio.
    \begin{center}
        \begin{tikzpicture}
            \begin{axis}[grafica, width=0.5\linewidth, height=4.5cm, ybar, bar width=5pt, xmin=23, xmax=27, ymin=0, ymax=95,
                xtick={23,24,25,26,27}, xlabel={Masa isotópica (u)}, ylabel={Abundancia (\%)}, grid=none,
                nodes near coords, nodes near coords style={font=\scriptsize, /pgf/number format/use comma}]
                \addplot[fill=ocre!70, draw=ocre] coordinates {(23.985,78.99) (24.986,10.00) (25.983,11.01)};
            \end{axis}
        \end{tikzpicture}
    \end{center}
    \begin{solu}
        \(\bar{A} = \dfrac{\num{23.985} \cdot \num{78.99} + \num{24.986} \cdot \num{10.00} + \num{25.983} \cdot \num{11.01}}{100} =
        \textcolor{red}{\qty{24.31}{\uma}}\).
    \end{solu}
\end{exercise}

\begin{exercise}
    \(\star\star\) El boro tiene dos isótopos, \ch{^{10}B} y \ch{^{11}B}, de masas \qty{10.013}{\uma} y \qty{11.009}{\uma}. Si
    la masa atómica del boro es \qty{10.81}{\uma}, calcula la abundancia natural de cada isótopo y representa su espectro de masas.
    \begin{solu}
        \(\num{10.81} = \dfrac{\num{10.013}\,x + \num{11.009}\,(100 - x)}{100}\Rightarrow\) \textcolor{red}{\ch{^{10}B}: 20,0\,\%}
        y \textcolor{red}{\ch{^{11}B}: 80,0\,\%}. El espectro tiene dos picos, el de 11 cuatro veces más alto que el de 10.
    \end{solu}
\end{exercise}

\begin{exercise}
    \(\star\star\) Calcula la masa atómica del plomo sabiendo que tiene cuatro isótopos: \ch{^{204}Pb} (\qty{203.97}{\uma};
    1,4\,\%), \ch{^{206}Pb} (\qty{205.97}{\uma}; 24,1\,\%), \ch{^{207}Pb} (\qty{206.98}{\uma}; 22,1\,\%) y \ch{^{208}Pb}
    (\qty{207.98}{\uma}).
    \begin{solu}
        Abundancia del \ch{^{208}Pb}: \(100 - \num{47.6} = \num{52.4}\)\,\%.
        \[ \bar{A} = \frac{\num{203.97} \cdot \num{1.4} + \num{205.97} \cdot \num{24.1} + \num{206.98} \cdot \num{22.1} + \num{207.98} \cdot \num{52.4}}{100}
        = \textcolor{red}{\qty{207.2}{\uma}} \]
    \end{solu}
\end{exercise}

\begin{exercise}
    \(\star\) El espectro de masas del litio da estos datos: \ch{^{6}Li}, \qty{6.015}{\uma} y 7,59\,\%; \ch{^{7}Li},
    \qty{7.016}{\uma} y 92,41\,\%. Calcula la masa atómica media del litio.
    \begin{solu}
        \(\bar{A} = \dfrac{\num{6.015} \cdot \num{7.59} + \num{7.016} \cdot \num{92.41}}{100} = \textcolor{red}{\qty{6.94}{\uma}}\).
    \end{solu}
\end{exercise}

\begin{exercise}
    \(\star\star\) El neón tiene tres isótopos estables: \ch{^{20}Ne} (\qty{19.992}{\uma}; 90,48\,\%), \ch{^{21}Ne}
    (\qty{20.994}{\uma}; 0,27\,\%) y \ch{^{22}Ne} (\qty{21.991}{\uma}; 9,25\,\%). Calcula su masa atómica. ¿Por qué el pico del
    \ch{^{21}Ne} es casi invisible en el espectro de masas?
    \begin{solu}
        \(\bar{A} = \textcolor{red}{\qty{20.18}{\uma}}\). \textcolor{red}{Porque su abundancia es solo del 0,27\,\%}: su pico es
        unas 335 veces más bajo que el del \ch{^{20}Ne}.
    \end{solu}
\end{exercise}

\begin{exercise}
    \(\star\star\star\) Representa el espectro de masas de estos elementos y calcula su masa atómica media: a) plata:
    \ch{^{107}Ag} (\qty{106.905}{\uma}; 51,84\,\%) y \ch{^{109}Ag} (\qty{108.905}{\uma}; 48,16\,\%); b) estroncio: \ch{^{84}Sr}
    (\qty{83.913}{\uma}; 0,56\,\%), \ch{^{86}Sr} (\qty{85.909}{\uma}; 9,86\,\%), \ch{^{87}Sr} (\qty{86.909}{\uma}; 7,00\,\%) y
    \ch{^{88}Sr} (\qty{87.906}{\uma}; 82,58\,\%).
    \begin{solu}
        a) \(\textcolor{red}{\qty{107.87}{\uma}}\): dos picos casi iguales. b) \(\textcolor{red}{\qty{87.62}{\uma}}\): un pico
        dominante en 88 y tres pequeños en 84, 86 y 87.
    \end{solu}
\end{exercise}

\subsection{Espectros atómicos}

\begin{exercise}
    \(\star\star\) A principios del siglo XX, Georges Claude inventó la lámpara de neón. a) Investiga qué colores pueden tener
    estas lámparas y de qué depende. b) El neón emite sobre todo luz rojo anaranjada, de unos \qty{640}{\nm}. ¿Cuál es su
    frecuencia? c) Si se sustituye el neón por argón, la luz es azul violácea (unos \qty{420}{\nm}). ¿Cuál es su frecuencia?
    \begin{solu}
        a) \textcolor{red}{El color depende del gas (cada uno tiene su espectro de emisión)}, y también del vidrio coloreado o de
        recubrimientos fluorescentes del tubo. b) \(\nu = \dfrac{\num{3.00e8}}{\num{640e-9}} = \textcolor{red}{\qty{4.69e14}{\hertz}}\).
        c) \(\textcolor{red}{\qty{7.14e14}{\hertz}}\).
    \end{solu}
\end{exercise}

\begin{exercise}
    \(\star\star\) La banda de radio de onda media (AM) va de 530 a \qty{1600}{\kilo\hertz}. a) ¿Entre qué longitudes de onda está?
    b) ¿Cuál es la longitud de onda de una emisora que emite a \qty{1008}{\kilo\hertz}?
    \begin{solu}
        a) \(\dfrac{\num{3.00e8}}{\num{530e3}} = \qty{566}{\m}\) y \(\dfrac{\num{3.00e8}}{\num{1600e3}} = \qty{188}{\m}\):
        \textcolor{red}{entre 188 y \qty{566}{\m}}. b) \(\textcolor{red}{\qty{298}{\m}}\).
    \end{solu}
\end{exercise}

\begin{exercise}
    \(\star\star\) Con la expresión de Bohr para la energía de los niveles del hidrógeno, calcula la longitud de onda de las cuatro
    rayas visibles de la serie de Balmer (saltos que terminan en \(n = 2\) desde \(n = 3, 4, 5\) y 6). Identifícalas en la
    figura~\ref{fig:ud3-espectros}.
    \begin{solu}
        \(\lambda = \dfrac{h c}{\num{2.18e-18} \left( \frac{1}{4} - \frac{1}{n^2} \right)}\): \(n = 3\): \textcolor{red}{\qty{657}{\nm}}
        (rojo); \(n = 4\): \textcolor{red}{\qty{486}{\nm}} (azul verdoso); \(n = 5\): \textcolor{red}{\qty{434}{\nm}} (azul
        violeta); \(n = 6\): \textcolor{red}{\qty{410}{\nm}} (violeta).
    \end{solu}
\end{exercise}

\begin{exercise}
    \(\star\star\) En el espectro de emisión del bario hay una raya verde de \qty{553.5}{\nm}. ¿Qué energía tiene cada fotón de esa
    luz? ¿Qué diferencia de energía hay entre los dos niveles del átomo de bario que producen esa raya?
    \begin{solu}
        \(E = \dfrac{h c}{\lambda} = \textcolor{red}{\qty{3.59e-19}{\joule}}\), que es \textcolor{red}{la misma diferencia de energía
        entre los dos niveles}.
    \end{solu}
\end{exercise}

\begin{exercise}
    \(\star\star\star\) Una radiación de \qty{680}{\nm} se ve de color rojo intenso. ¿Qué energía tiene uno de sus fotones, en julios
    y en electronvoltios? ¿Y un mol de fotones?
    \begin{solu}
        \(E = \dfrac{\num{6.626e-34} \cdot \num{3.00e8}}{\num{680e-9}} = \textcolor{red}{\qty{2.92e-19}{\joule}}\);
        \(\dfrac{\num{2.92e-19}}{\num{1.602e-19}} = \textcolor{red}{\qty{1.82}{\electronvolt}}\); un mol:
        \(\textcolor{red}{\qty{1.76e5}{\joule}}\).
    \end{solu}
\end{exercise}

\begin{exercise}
    \(\star\) La energía de ionización de un átomo de litio es \qty{8.637e-19}{\joule}. ¿Cuál es la longitud de onda máxima de la
    radiación capaz de ionizarlo? ¿En qué zona del espectro está?
    \begin{solu}
        \(\lambda = \dfrac{h c}{E} = \dfrac{\num{1.988e-25}}{\num{8.637e-19}} = \textcolor{red}{\qty{2.30e-7}{\m} = \qty{230}{\nm}}\),
        en el \textcolor{red}{ultravioleta}. Una radiación de longitud de onda mayor no tiene fotones con energía suficiente.
    \end{solu}
\end{exercise}

\begin{exercise}
    \(\star\star\) Calcula la longitud de onda de la raya de la serie de Paschen del hidrógeno que corresponde al salto de
    \(n = 7\) a \(n = 3\).
    \begin{solu}
        \(\dfrac{1}{\lambda} = \num{1.097e7} \left( \dfrac{1}{9} - \dfrac{1}{49} \right) = \qty{9.95e5}{\per\m}\):
        \(\lambda = \textcolor{red}{\qty{1.005e-6}{\m} = \qty{1005}{\nm}}\) (infrarrojo).
    \end{solu}
\end{exercise}

\begin{exercise}
    \(\star\) ¿Cuántas veces es más energético un fotón ultravioleta de \qty{250}{\nm} que uno infrarrojo de \qty{2}{\micro\m}?
    \begin{solu}
        Como \(E = hc/\lambda\), la relación es la inversa de las longitudes de onda: \(\dfrac{2000}{250} = \textcolor{red}{8}\)
        veces (\qty{7.95e-19}{\joule} frente a \qty{9.94e-20}{\joule}).
    \end{solu}
\end{exercise}

\subsection{Números cuánticos}

\begin{exercise}
    \(\star\) El estado de un electrón queda descrito por cuatro números cuánticos. ¿Qué nos indica cada uno y qué valores pueden
    tomar?
    \begin{solu}
        \(\textcolor{red}{n}\) (\(1, 2, 3\dots\)): el nivel de energía y el tamaño del orbital. \(\textcolor{red}{l}\) (de 0 a
        \(n-1\)): el subnivel y la forma del orbital. \(\textcolor{red}{m}\) (de \(-l\) a \(+l\)): la orientación del orbital.
        \(\textcolor{red}{m_s}\) (\(\pm\frac{1}{2}\)): el sentido del espín del electrón.
    \end{solu}
\end{exercise}

\begin{exercise}
    \(\star\star\) Si \(n = 2\), ¿qué valores pueden tomar los demás números cuánticos? ¿Cuántos electrones caben en ese nivel?
    \begin{solu}
        \(l = 0\): \(m = 0\); \(l = 1\): \(m = -1, 0, +1\); en todos los casos \(m_s = \pm\frac{1}{2}\). Hay 4 orbitales:
        \textcolor{red}{8 electrones}.
    \end{solu}
\end{exercise}

\begin{exercise}
    \(\star\star\) Indica en qué orbital se encuentra el electrón descrito por cada conjunto de números cuánticos:
    \((2, 0, 0, +\frac{1}{2})\); \((3, 2, -1, -\frac{1}{2})\); \((4, 1, 0, -\frac{1}{2})\); \((5, 3, -2, +\frac{1}{2})\).
    \begin{solu}
        \textcolor{red}{2s; 3d; 4p; 5f.}
    \end{solu}
\end{exercise}

\begin{exercise}
    \(\star\) En el subnivel 4s hay dos electrones. ¿Cuáles son sus números cuánticos?
    \begin{solu}
        \(\textcolor{red}{(4, 0, 0, +\frac{1}{2})}\) y \(\textcolor{red}{(4, 0, 0, -\frac{1}{2})}\): por el principio de Pauli,
        solo se diferencian en el espín.
    \end{solu}
\end{exercise}

\begin{exercise}
    \(\star\star\) ¿Cuáles de estas combinaciones de números cuánticos están permitidas? Explica por qué.
    a) \((4, 2, -3, -\frac{1}{2})\); b) \((3, 3, 1, +\frac{1}{2})\); c) \((2, 1, 0, -\frac{1}{2})\); d) \((1, 0, 1, +\frac{1}{2})\).
    \begin{solu}
        a) \textcolor{red}{No}: con \(l = 2\), \(m\) va de \(-2\) a \(+2\). b) \textcolor{red}{No}: \(l\) debe ser menor que \(n\).
        c) \textcolor{red}{Sí} (orbital 2p). d) \textcolor{red}{No}: con \(l = 0\), \(m\) solo puede ser 0.
    \end{solu}
\end{exercise}

\begin{exercise}
    \(\star\star\) Indica el número máximo de electrones que puede haber en cada caso y en qué orbitales están:
    a) \(n = 3\) y \(m_s = -\frac{1}{2}\); b) \(n = 4\), \(l = 2\) y \(m = -1\); c) \(n = 5\) y \(l = 1\).
    \begin{solu}
        a) El nivel 3 tiene 9 orbitales (3s, 3p, 3d) y en cada uno hay un electrón con \(m_s = -\frac{1}{2}\):
        \textcolor{red}{9 electrones}. b) Un orbital 4d: \textcolor{red}{2 electrones}. c) Los tres orbitales 5p:
        \textcolor{red}{6 electrones}.
    \end{solu}
\end{exercise}

\begin{exercise}
    \(\star\star\) ¿Son verdaderas o falsas estas afirmaciones? a) Un electrón de un orbital 3s puede tener los números cuánticos
    \((3, 1, 0, -\frac{1}{2})\). b) La combinación \((3, 1, -1, -\frac{1}{2})\) está permitida. c) En el nivel \(n = 2\) caben como
    máximo 10 electrones.
    \begin{solu}
        a) \textcolor{red}{Falsa}: en un orbital s, \(l = 0\). b) \textcolor{red}{Verdadera} (es un electrón 3p). c)
        \textcolor{red}{Falsa}: caben \(2n^2 = 8\).
    \end{solu}
\end{exercise}

\begin{exercise}
    \(\star\star\) Escribe una combinación válida de números cuánticos para un electrón de cada uno de estos subniveles: 2s, 4d y 5p.
    \begin{solu}
        Por ejemplo: \(\textcolor{red}{(2, 0, 0, +\frac{1}{2})}\); \(\textcolor{red}{(4, 2, 1, -\frac{1}{2})}\);
        \(\textcolor{red}{(5, 1, -1, +\frac{1}{2})}\).
    \end{solu}
\end{exercise}

\begin{exercise}
    \(\star\star\) Escribe una combinación válida de números cuánticos para: a) un orbital 3p; b) un electrón 3p. ¿Por qué no hace
    falta el mismo número de datos en los dos casos?
    \begin{solu}
        a) \(\textcolor{red}{(n, l, m) = (3, 1, 0)}\), por ejemplo. b) \(\textcolor{red}{(3, 1, 0, +\frac{1}{2})}\).
        \textcolor{red}{Un orbital queda definido por \(n\), \(l\) y \(m\); un electrón necesita además su espín}, porque en cada
        orbital puede haber dos.
    \end{solu}
\end{exercise}

\subsection{Configuración electrónica}

\begin{exercise}
    \(\star\) A partir de la configuración electrónica del cloro, \conf{1s^2\,2s^2\,2p^6\,3s^2\,3p^5}, indica: a) su número atómico
    y su periodo; b) el ion más probable; c) su grupo.
    \begin{solu}
        a) \(\textcolor{red}{Z = 17}\), \textcolor{red}{periodo 3}. b) \textcolor{red}{\ch{Cl^-}}: con un electrón más completa su
        capa de valencia. c) \textcolor{red}{Grupo 17} (halógenos).
    \end{solu}
\end{exercise}

\begin{exercise}
    \(\star\) A partir de la configuración electrónica del azufre, \conf{1s^2\,2s^2\,2p^6\,3s^2\,3p^4}, indica: a) su número
    atómico y su periodo; b) el ion más probable; c) su grupo.
    \begin{solu}
        a) \(\textcolor{red}{Z = 16}\), \textcolor{red}{periodo 3}. b) \textcolor{red}{\ch{S^2-}}. c) \textcolor{red}{Grupo 16}.
    \end{solu}
\end{exercise}

\begin{exercise}
    \(\star\) Para escribir la configuración electrónica de un átomo hay que seguir tres reglas. ¿Cuáles son? Enúncialas.
    \begin{solu}
        \textcolor{red}{Principio de Aufbau} (los electrones ocupan los subniveles de menor a mayor energía),
        \textcolor{red}{principio de exclusión de Pauli} (no hay dos electrones con los cuatro números cuánticos iguales: dos como
        máximo por orbital) y \textcolor{red}{regla de Hund} (en un subnivel, máximo número de electrones desapareados).
    \end{solu}
\end{exercise}

\begin{exercise}
    \(\star\star\) ¿Cuáles de estas configuraciones corresponden a átomos con propiedades parecidas? a) \conf{1s^2\,2s^2\,2p^6\,3s^1};
    b) \conf{[Ar]\,4s^2\,3d^{10}\,4p^5}; c) \conf{1s^2\,2s^2\,2p^6\,3s^2\,3p^3}; d) \conf{1s^2\,2s^1}; e) \conf{1s^2\,2s^2\,2p^3};
    f) \conf{1s^2\,2s^2\,2p^6\,3s^2\,3p^5}; g) \conf{[Ar]\,4s^1}; h) \conf{[Ar]\,4s^2\,3d^{10}\,4p^3}.
    \begin{solu}
        \textcolor{red}{a, d y g} (\conf{ns^1}, grupo 1: Na, Li, K); \textcolor{red}{b y f} (\conf{ns^2\,np^5}, grupo 17: Br, Cl);
        \textcolor{red}{c, e y h} (\conf{ns^2\,np^3}, grupo 15: P, N, As).
    \end{solu}
\end{exercise}

\begin{exercise}
    \(\star\star\) Escribe la configuración electrónica de estos iones e indica cuáles tienen electrones desapareados: \ch{Zn^2+}
    (\(Z = 30\)), \ch{Br^-} (\(Z = 35\)), \ch{Ga^3+} (\(Z = 31\)), \ch{Se^2-} (\(Z = 34\)), \ch{Fe^3+} (\(Z = 26\)) y \ch{Mg^2+}
    (\(Z = 12\)).
    \begin{solu}
        \ch{Zn^2+}: \conf{[Ar]\,3d^{10}}; \ch{Br^-}: \conf{[Ar]\,4s^2\,3d^{10}\,4p^6}; \ch{Ga^3+}: \conf{[Ar]\,3d^{10}};
        \ch{Se^2-}: \conf{[Ar]\,4s^2\,3d^{10}\,4p^6}; \ch{Fe^3+}: \conf{[Ar]\,3d^5}; \ch{Mg^2+}: \conf{1s^2\,2s^2\,2p^6}.
        \textcolor{red}{Solo el \ch{Fe^3+} tiene electrones desapareados} (5, uno en cada orbital 3d); los demás tienen todos los
        electrones apareados.
    \end{solu}
\end{exercise}

\begin{exercise}
    \(\star\) Para las configuraciones de la capa de valencia \conf{ns^2\,np^1}, \conf{ns^2\,np^4} y \conf{ns^2\,np^5}: a) indica
    el grupo de cada una; b) nombra un elemento de cada grupo; c) ¿cuál será el ion más estable que formarán? Razónalo.
    \begin{solu}
        a) \textcolor{red}{Grupos 13, 16 y 17.} b) Por ejemplo, \textcolor{red}{aluminio, oxígeno y flúor}. c)
        \textcolor{red}{\ch{X^3+}, \ch{X^2-} y \ch{X^-}}: el del grupo 13 pierde sus tres electrones de valencia y los de los grupos
        16 y 17 ganan los electrones que les faltan para completar la capa (configuración de gas noble).
    \end{solu}
\end{exercise}

\begin{exercise}
    \(\star\) Razona si estas configuraciones corresponden a un estado fundamental, a un estado excitado o si son imposibles:\newline
    a) \conf{1s^2\,2s^2\,2p^6\,3s^2\,3p^4\,4s^1}; b) \conf{1s^2\,2s^2\,2p^6\,2d^1}; c) \conf{1s^2\,2s^2\,2p^6\,3s^2\,3p^6\,4s^1}.
    \begin{solu}
        a) \textcolor{red}{Excitado}: el 3p no está completo y ya hay un electrón en 4s. b) \textcolor{red}{Imposible}: no existe
        el subnivel 2d (con \(n = 2\), \(l\) solo puede ser 0 o 1). c) \textcolor{red}{Fundamental} (es el potasio).
    \end{solu}
\end{exercise}

\begin{exercise}
    \(\star\star\) Para el elemento de número atómico 33: a) escribe su configuración electrónica; b) ¿a qué grupo y periodo
    pertenece?; c) ¿qué valores tienen los números cuánticos de su último electrón?
    \begin{solu}
        a) \(\textcolor{red}{\conf{1s^2\,2s^2\,2p^6\,3s^2\,3p^6\,4s^2\,3d^{10}\,4p^3}}\). b) \textcolor{red}{Grupo 15, periodo 4}
        (arsénico). c) El tercer electrón 4p ocupa, según la regla de Hund, el último orbital p con espín \(+\frac{1}{2}\):
        \(\textcolor{red}{(4, 1, +1, +\frac{1}{2})}\).
    \end{solu}
\end{exercise}

\subsection{Propiedades periódicas}

\begin{exercise}
    \(\star\star\) ¿Qué tendrá mayor radio, un átomo en estado fundamental o su catión? ¿Y si en lugar de un catión fuera un anión?
    \begin{solu}
        \textcolor{red}{El átomo es mayor que su catión}: al perder electrones disminuye la repulsión entre ellos (y a menudo
        desaparece la última capa). \textcolor{red}{El anión es mayor que el átomo}: el electrón ganado aumenta la repulsión y la
        nube se expande.
    \end{solu}
\end{exercise}

\begin{exercise}
    \(\star\) ¿Por qué no se suelen dar valores del radio atómico de los gases nobles?
    \begin{solu}
        \textcolor{red}{Porque el radio atómico se mide a partir de la distancia entre núcleos de átomos enlazados, y los gases
        nobles prácticamente no forman enlaces.}
    \end{solu}
\end{exercise}

\begin{exercise}
    \(\star\) Ordena de menor a mayor radio atómico: N, P, As, Sb y Bi. Justifica la respuesta.
    \begin{solu}
        Son del grupo 15; el radio aumenta al bajar porque cada uno tiene una capa más:
        \textcolor{red}{N \(<\) P \(<\) As \(<\) Sb \(<\) Bi}.
    \end{solu}
\end{exercise}

\begin{exercise}
    \(\star\star\) Razona cuál tiene mayor radio en cada pareja: a) Cl y \ch{Cl^-}; b) K y \ch{K^+}.
    \begin{solu}
        a) \textcolor{red}{\ch{Cl^-}}: tiene un electrón más con los mismos protones. b) \textcolor{red}{K}: el \ch{K^+} ha perdido
        la capa 4.
    \end{solu}
\end{exercise}

\begin{exercise}
    \(\star\star\) Para los elementos de números atómicos 8 y 11: a) escribe su configuración electrónica; b) indica el ion más
    estable de cada uno y su configuración; c) ¿cuál de los dos iones tendrá mayor radio iónico?
    \begin{solu}
        a) \(Z = 8\): \conf{1s^2\,2s^2\,2p^4} (O); \(Z = 11\): \conf{1s^2\,2s^2\,2p^6\,3s^1} (Na). b) \textcolor{red}{\ch{O^2-} y
        \ch{Na^+}}, los dos con configuración \conf{1s^2\,2s^2\,2p^6}. c) Tienen los mismos electrones (10), pero el oxígeno solo
        tiene 8 protones y el sodio 11: \textcolor{red}{el \ch{O^2-} es mayor}.
    \end{solu}
\end{exercise}

\begin{exercise}
    \(\star\star\) Ordena estos elementos de menor a mayor primera energía de ionización: K, Ca, Fe y Cl. Si tuviéramos en cuenta la
    segunda energía de ionización, ¿estarían en el mismo orden?
    \begin{solu}
        Los cuatro están en el periodo 3 o 4 y la energía de ionización aumenta hacia la derecha y hacia arriba:
        \textcolor{red}{K \(<\) Ca \(<\) Fe \(<\) Cl} (419, 590, 762 y 1251 kJ/mol).
        \textcolor{red}{No}: el \ch{K^+} tiene la configuración del argón y arrancarle otro electrón cuesta muchísimo, así que el K
        pasa a ser el de mayor segunda energía de ionización: \(\mathrm{Ca} < \mathrm{Fe} < \mathrm{Cl} < \mathrm{K}\)
        (1145, 1562, 2298 y 3052 kJ/mol).
    \end{solu}
\end{exercise}

\begin{exercise}
    \(\star\star\) Las tres primeras energías de ionización del magnesio son 738, 1451 y \qty{7733}{\kilo\joule\per\mol}.
    a) ¿Por qué la segunda es aproximadamente el doble de la primera? b) ¿Por qué la tercera es tanto mayor que la segunda?
    \begin{solu}
        a) \textcolor{red}{Los dos primeros electrones salen del mismo subnivel 3s}, pero el segundo se arranca a un ion positivo,
        que lo atrae más. b) \textcolor{red}{El tercero hay que sacarlo del nivel 2, completo y mucho más cercano al núcleo}.
    \end{solu}
\end{exercise}

\begin{exercise}
    \(\star\star\) A partir de estas configuraciones electrónicas fundamentales:\newline a) \conf{1s^2\,2s^2\,2p^6\,3s^2\,3p^2};
    b) \conf{1s^2\,2s^2\,2p^6\,3s^2\,3p^5}; c) \conf{1s^2\,2s^2\,2p^6\,3s^2\,3p^6\,4s^1}; d) \conf{1s^2\,2s^2\,2p^6\,3s^2\,3p^6\,4s^2},
    ordénalos de mayor a menor primera energía de ionización y de mayor a menor segunda energía de ionización.
    \begin{solu}
        Son Si, Cl, K y Ca. Primera: \(\textcolor{red}{\mathrm{b} > \mathrm{a} > \mathrm{d} > \mathrm{c}}\) (Cl \(>\) Si \(>\)
        Ca \(>\) K). Segunda: \(\textcolor{red}{\mathrm{c} > \mathrm{b} > \mathrm{a} > \mathrm{d}}\): el \ch{K^+} tiene la
        configuración de gas noble, y el \ch{Ca^+} pierde fácilmente su último electrón 4s para quedar también como el argón.
    \end{solu}
\end{exercise}

\begin{exercise}
    \(\star\star\) Para la pareja de átomos Mg y Cl: a) ¿cuál tiene mayor radio?; b) ¿cuál tiene mayor energía de ionización?;
    c) ¿cuál es más electronegativo?
    \begin{solu}
        Los dos están en el periodo 3, el Mg en el grupo 2 y el Cl en el 17. a) \textcolor{red}{Mg}. b) \textcolor{red}{Cl}.
        c) \textcolor{red}{Cl}.
    \end{solu}
\end{exercise}

\begin{exercise}
    \(\star\star\) Para la pareja de átomos Rb y Br: a) ¿cuál tiene menor radio?; b) ¿cuál tiene mayor energía de ionización?;
    c) ¿cuál tiene menor electronegatividad?
    \begin{solu}
        El Rb está en el grupo 1 y el periodo 5; el Br, en el grupo 17 y el periodo 4. a) \textcolor{red}{Br}. b)
        \textcolor{red}{Br}. c) \textcolor{red}{Rb}.
    \end{solu}
\end{exercise}

\begin{exercise}
    \(\star\) ¿Podríamos saber si un elemento es metálico o no metálico a partir del valor de su electronegatividad? Razona la
    respuesta.
    \begin{solu}
        \textcolor{red}{Sí, de forma aproximada}: los metales tienen electronegatividades bajas (en general, menores de 2), porque
        tienden a perder electrones; los no metales, altas, porque tienden a captarlos.
    \end{solu}
\end{exercise}

%========================================================================================
\section{Ejercicios de ampliación}\label{sec:ud3-ampliacion}
%========================================================================================

\subsection{Estructura atómica}

\begin{exercise}
    \(\star\star\) Haz un trabajo de investigación sobre el experimento de la gota de aceite de Millikan. ¿Por qué, muchos años
    después, se acusó a Millikan de haber seleccionado sus datos? ¿Estaba justificada la polémica?
    \begin{solu}
        \textcolor{red}{Respuesta abierta.} \textcolor{blue}{Millikan medía la velocidad de gotitas de aceite cargadas entre dos
        placas y calculaba su carga: todas eran múltiplos de \qty{1.6e-19}{\coulomb}. Sus cuadernos mostraron que había descartado
        algunas medidas que no publicó. Se discutió si era un fraude; hoy se piensa que descartó sobre todo medidas con problemas
        experimentales y que su resultado era correcto, aunque su forma de informar no fue del todo transparente.}
    \end{solu}
\end{exercise}

\begin{exercise}
    \(\star\star\) El potasio se usa en algunas células fotoeléctricas porque basta una energía pequeña, \qty{2.30}{\electronvolt},
    para arrancarle electrones de la superficie del metal. ¿Cuál es la longitud de onda máxima de la luz capaz de arrancar
    electrones a una placa de potasio? ¿Funcionaría la célula con luz roja?
    \begin{solu}
        \(E = \num{2.30} \cdot \num{1.602e-19} = \qty{3.68e-19}{\joule}\); \(\lambda = \dfrac{hc}{E} = \textcolor{red}{\qty{5.39e-7}{\m}
        = \qty{539}{\nm}}\). \textcolor{red}{No}: la luz roja tiene mayor longitud de onda y sus fotones no tienen energía suficiente,
        por intensa que sea.
    \end{solu}
\end{exercise}

\begin{exercise}
    \(\star\star\) ¿Cuál de estas radiaciones electromagnéticas tiene mayor longitud de onda? a) \qty{450}{\nm}; b) \qty{500}{\mega\hertz};
    c) \qty{3000}{\angstrom}; d) \qty{25000}{\per\cm}.
    \begin{solu}
        a) \qty{4.5e-7}{\m}; b) \(\lambda = \dfrac{\num{3.00e8}}{\num{5.00e8}} = \qty{0.6}{\m}\); c) \qty{3e-7}{\m}; d)
        \(\lambda = 1/\tilde{\nu} = \qty{4e-5}{\cm} = \qty{4e-7}{\m}\). \textcolor{red}{La b} (es una onda de radio).
    \end{solu}
\end{exercise}

\begin{exercise}
    \(\star\star\star\) Al iluminar un átomo de magnesio en estado gaseoso con radiación de \qty{162.2}{\nm} se le arranca un
    electrón. a) ¿Cuál es la primera energía de ionización del magnesio, en \unit{\kilo\joule\per\mol}? b) Para arrancarle un
    segundo electrón se necesita radiación de \qty{82.5}{\nm}. ¿Cuál es la segunda energía de ionización? Compara los valores.
    \begin{solu}
        a) \(E = \dfrac{h c}{\lambda} = \qty{1.226e-18}{\joule}\) por átomo; \(\cdot N_A = \textcolor{red}{\qty{738}{\kilo\joule\per\mol}}\).
        b) \(\qty{2.409e-18}{\joule}\); \(\textcolor{red}{\qty{1451}{\kilo\joule\per\mol}}\). \textcolor{red}{La segunda es casi
        el doble}: el electrón se arranca a un ion \ch{Mg^+}, con más carga positiva que lo atrae.
    \end{solu}
\end{exercise}

\begin{exercise}
    \(\star\star\star\) Explica cuáles de estos conjuntos de números cuánticos son incorrectos: a) \((4, 3, -2, +\frac{1}{2})\);
    b) \((2, 1, -1, 0)\); c) \((3, 0, 1, -\frac{1}{2})\); d) \((1, 1, 0, +\frac{1}{2})\).
    \begin{solu}
        a) Correcto (electrón 4f). \textcolor{red}{b) Incorrecto}: \(m_s\) no puede ser 0. \textcolor{red}{c) Incorrecto}: con
        \(l = 0\), \(m = 0\). \textcolor{red}{d) Incorrecto}: con \(n = 1\), \(l\) solo puede ser 0.
    \end{solu}
\end{exercise}

\begin{exercise}
    \(\star\star\star\) ¿Cuántos electrones con números cuánticos distintos puede haber en un subnivel con \(n = 4\) y \(l = 2\)?
    Escribe los conjuntos de números cuánticos de todos ellos.
    \begin{solu}
        Es el subnivel 4d: 5 orbitales (\(m = -2, -1, 0, +1, +2\)), cada uno con dos espines: \textcolor{red}{10 electrones}:
        \((4, 2, m, \pm\frac{1}{2})\) con \(m = -2, -1, 0, +1, +2\).
    \end{solu}
\end{exercise}

\begin{exercise}
    \(\star\star\star\) El electrón de valencia más externo de un átomo en estado fundamental tiene los números cuánticos
    \(n = 4\), \(l = 1\), \(m = 0\), \(m_s = -\frac{1}{2}\). Suponiendo que los orbitales se llenan de \(m = -l\) a \(m = +l\) y
    primero con espín \(+\frac{1}{2}\), ¿de qué elemento se trata?
    \begin{solu}
        Es un electrón 4p. Con esa convención, los tres primeros electrones 4p tienen \(m_s = +\frac{1}{2}\) (\(m = -1, 0, +1\)); el
        cuarto, \((4, 1, -1, -\frac{1}{2})\), y el quinto, \((4, 1, 0, -\frac{1}{2})\). La configuración termina en \conf{4p^5}:
        \conf{[Ar]\,4s^2\,3d^{10}\,4p^5}, \(Z = 35\): \textcolor{red}{bromo}.
    \end{solu}
\end{exercise}

\begin{exercise}
    \(\star\star\star\) El electrón de un átomo de hidrógeno en su estado fundamental absorbe un fotón de \qty{97.3}{\nm}. ¿A qué
    nivel salta? ¿Qué radiaciones podrá emitir después al volver al estado fundamental?
    \begin{solu}
        \(E = \dfrac{hc}{\lambda} = \qty{2.04e-18}{\joule} = \num{2.18e-18} \left( 1 - \dfrac{1}{n^2} \right) \Rightarrow
        n^2 = 16\): \textcolor{red}{salta al nivel \(n = 4\)}. Al volver puede hacerlo directamente (\(4 \rightarrow 1\), en el
        ultravioleta) o por etapas: \textcolor{red}{\(4 \rightarrow 3\) y \(3 \rightarrow 2\) (infrarrojo y visible),
        \(4 \rightarrow 2\) (visible, \qty{486}{\nm}), \(3 \rightarrow 1\) y \(2 \rightarrow 1\) (ultravioleta)}: seis rayas
        posibles.
    \end{solu}
\end{exercise}

\subsection{Tabla periódica}

\begin{exercise}
    \(\star\star\star\) Para el elemento de configuración electrónica
    \[ \conf{1s^2\,2s^2\,2p^6\,3s^2\,3p^6\,4s^2\,3d^{10}\,4p^4} \]
    a) ¿en qué periodo y grupo está?; b) ¿qué elemento es?; c) nombra dos elementos de su mismo grupo; d) escribe los números
    cuánticos de su electrón diferenciador; e) representa los electrones de valencia en un diagrama de cajas. ¿Tiene electrones
    desapareados?
    \begin{solu}
        a) \textcolor{red}{Periodo 4, grupo 16.} b) \(Z = 34\): \textcolor{red}{selenio}. c) \textcolor{red}{Oxígeno y azufre}
        (o teluro). d) El cuarto electrón 4p se aparea en el primer orbital: \(\textcolor{red}{(4, 1, -1, -\frac{1}{2})}\).
        e) 4s \orbcaja{\uparrow\downarrow} 4p \orbcaja{\uparrow\downarrow}\orbcaja{\uparrow}\orbcaja{\uparrow}:
        \textcolor{red}{dos electrones desapareados}.
    \end{solu}
\end{exercise}

\begin{exercise}
    \(\star\star\) ¿Cuál de estas configuraciones electrónicas corresponde a un átomo en estado fundamental, cuál a un estado
    excitado y cuál es imposible? Razona la respuesta. a) \conf{1s^2\,2s^3\,2p^6\,3s^1}; b) \conf{[Ar]\,4s^2\,3d^{10}\,4p^3};
    c) \conf{1s^2\,2s^2\,2p^5\,3s^1}.
    \begin{solu}
        a) \textcolor{red}{Imposible}: en el orbital 2s solo caben dos electrones (Pauli). b) \textcolor{red}{Fundamental} (arsénico).
        c) \textcolor{red}{Excitado}: un electrón ha pasado del 2p al 3s (es el neón excitado).
    \end{solu}
\end{exercise}

\begin{exercise}
    \(\star\star\) Para el átomo de número atómico 53, razona si son verdaderas o falsas estas afirmaciones: a) está en el grupo 16;
    b) tiene los orbitales 4d completos; c) es un metal de transición; d) si capta un electrón, su anión tiene la configuración del
    gas noble siguiente.
    \begin{solu}
        Su configuración es \conf{[Kr]\,5s^2\,4d^{10}\,5p^5} (yodo). a) \textcolor{red}{Falsa}: grupo 17. b) \textcolor{red}{Verdadera}.
        c) \textcolor{red}{Falsa}: es un halógeno (no metal, bloque p). d) \textcolor{red}{Verdadera}: \ch{I^-} tiene la
        configuración del xenón.
    \end{solu}
\end{exercise}

\begin{exercise}
    \(\star\star\) Para las configuraciones de dos átomos neutros, A: \conf{1s^2\,2s^2\,2p^6\,3s^2\,3p^6\,4s^1} y
    B: \conf{1s^2\,2s^2\,2p^6\,3s^2\,3p^6\,5s^1}, razona si son ciertas estas afirmaciones: a) A es un átomo de potasio; b) A y B son
    átomos de elementos distintos; c) al pasar de B a A se desprende energía; d) B es un metal alcalinotérreo; e) cuesta más arrancar
    un electrón de A que de B.
    \begin{solu}
        Los dos tienen 19 electrones. a) \textcolor{red}{Cierta.} b) \textcolor{red}{Falsa}: es el mismo elemento; B es el potasio en
        un estado excitado. c) \textcolor{red}{Cierta}: el electrón baja de nivel y emite un fotón. d) \textcolor{red}{Falsa}: es un
        alcalino excitado. e) \textcolor{red}{Cierta}: en B el electrón está más lejos del núcleo y menos atraído.
    \end{solu}
\end{exercise}

\begin{exercise}
    \(\star\star\star\) De cada una de estas parejas de elementos: C y O; Li y K; Al y Cl; Ca y Ba, indica cuál tiene: a) mayor radio
    atómico; b) menor energía de ionización; c) mayor electronegatividad; d) menor carácter metálico.
    \begin{solu}
        \begin{tabular}{@{}lcccc@{}}
            \toprule
            & \textbf{C/O} & \textbf{Li/K} & \textbf{Al/Cl} & \textbf{Ca/Ba} \\
            \midrule
            Mayor radio & C & K & Al & Ba \\[3pt]
            Menor energía de ionización & C & K & Al & Ba \\[3pt]
            Mayor electronegatividad & O & Li & Cl & Ca \\[3pt]
            Menor carácter metálico & O & Li & Cl & Ca \\
            \bottomrule
        \end{tabular}

        \textcolor{red}{C/O y Al/Cl son del mismo periodo} (manda la posición a la derecha); \textcolor{red}{Li/K y Ca/Ba, del mismo
        grupo} (manda la posición hacia abajo).
    \end{solu}
\end{exercise}

\begin{exercise}
    \(\star\star\star\) El cloro natural tiene una masa atómica de \qty{35.45}{\uma} y está formado por \ch{^{35}Cl}
    (\qty{34.969}{\uma}) y \ch{^{37}Cl} (\qty{36.966}{\uma}). a) Calcula la abundancia de cada isótopo. b) Las moléculas de \ch{Cl2}
    pueden tener masa 70, 72 o 74 u. ¿En qué proporción aparecerá cada una en el espectro de masas del cloro gaseoso? (La
    probabilidad de que los dos átomos sean \ch{^{35}Cl} es el cuadrado de la abundancia de ese isótopo.)
    \begin{solu}
        a) \(\num{35.45} = \dfrac{\num{34.969}\,x + \num{36.966}\,(100 - x)}{100} \Rightarrow\) \textcolor{red}{\ch{^{35}Cl}: 75,9\,\%;
        \ch{^{37}Cl}: 24,1\,\%}. b) \(p = \num{0.759}\), \(q = \num{0.241}\): masa 70, \(p^2 =\) \textcolor{red}{57,6\,\%}; masa 72
        (un átomo de cada), \(2pq =\) \textcolor{red}{36,6\,\%}; masa 74, \(q^2 =\) \textcolor{red}{5,8\,\%}.
    \end{solu}
\end{exercise}
